% =============================================================================
% Referee-revised patch: §5 Toy anomaly --- polynomial-weight closure
%   sec:toy-anomaly-and-transverse-obstruction
%
% Purpose: replace the impossible uniform-in-D_T extension by a rigorous
% polynomial-weight retained subdomain.  This patch corrects two issues in the
% previous draft:
%   (i)  the retained set is signed via q(m)|s| in D, matching the density
%        proof that counts both signs of s;
%   (ii) the density estimate is stated using q_T^- and q_T^+, avoiding the
%        false assertion q(m) <= q(T) on a symmetric theta window.
% =============================================================================

% =============================================================================
% (1) Insert this subsection BETWEEN subsec:toy-anomaly-functional and
%     subsec:toy-anomaly-visibility.
% =============================================================================

\subsection{Toy retained subdomain and polynomial-weight density}
\label{subsec:toy-retained-subdomain}

\begin{definition}[Toy retained subdomain]
\label{def:toy-retained-subdomain}
Fix a compact interval \(D=[d_-,d_+]\subset(0,\infty)\) avoiding the
isolated zeros of \(F_{\infty}\) of
Lemma~\ref{lem:toy-anomaly-leading}.  The signed toy retained subdomain
at height \(T\) is
\begin{equation}
\label{eq:toy-retained-subdomain}
        \mathcal D_T^\toy(D)
        :=\bigl\{(m,s)\in\mathcal D_T^\times:
              q(m)\,|s|\in D\bigr\}.
\end{equation}
Thus the positive and negative separation branches are both retained.
The comparison estimates below are stated for the absolute scalar
functional \(|\mathcal A_\toy|\), so retaining both signs does not change
the toy lower bound; see Lemma~\ref{lem:toy-signed-separation-invariance}.
\end{definition}

\begin{lemma}[Polynomial-weight density of the toy retained subdomain]
\label{lem:toy-retained-density}
Let
\[
        q_T^-:=\inf_{m\in I_T}q(m),
        \qquad
        q_T^+:=\sup_{m\in I_T}q(m).
\]
Assume \(q_T^->0\) and
\begin{equation}
\label{eq:toy-density-height-hypothesis}
        \frac{d_+}{q_T^-}\le \frac{|I_T|}{4}.
\end{equation}
Then
\begin{equation}
\label{eq:toy-retained-density-exact}
        \frac{|\mathcal D_T^\toy(D)|}{|\mathcal D_T|}
        \ge
        \frac{d_+-d_-}{q_T^+\,|I_T|}.
\end{equation}
Consequently, if the retained windows satisfy
\begin{equation}
\label{eq:toy-window-comparability}
        |I_T|\le C_I Q(T)^{-1},
        \qquad
        q_T^+\le C_q q(T)
\end{equation}
with absolute constants \(C_I,C_q\), then
\begin{equation}
\label{eq:toy-retained-density}
        \frac{|\mathcal D_T^\toy(D)|}{|\mathcal D_T|}
        \ge
        \frac{d_+-d_-}{C_I C_q}\,\frac{Q(T)}{q(T)}.
\end{equation}
The height hypothesis \eqref{eq:toy-density-height-hypothesis} is
implied for all sufficiently large \(T\) in any scale regime with
\(|I_T|\ge c_I Q(T)^{-1}\), \(q_T^-\ge c_q q(T)\), and
\(Q(T)=o(q(T))\), with constants depending only on \(D,c_I,c_q\).
\end{lemma}

\begin{proof}
The finite packet domain is
\[
    \mathcal D_T
    =\bigl\{(m,s): |m-T|+|s|/2\le |I_T|/2\bigr\},
\]
a rotated square of area \(|\mathcal D_T|=|I_T|^2\).  Restrict to the
central slab
\[
        S_T:=\{m: |m-T|\le |I_T|/4\}.
\]
For \(m\in S_T\), the \(\mathcal D_T\)-constraint permits
\(|s|\le |I_T|/2-|m-T|\), hence at least \(|s|\le |I_T|/4\).  By
\eqref{eq:toy-density-height-hypothesis},
\[
        \frac{d_+}{q(m)}\le \frac{d_+}{q_T^-}\le \frac{|I_T|}{4},
\]
so both intervals
\[
        \left[\frac{d_-}{q(m)},\frac{d_+}{q(m)}\right],
        \qquad
        -\left[\frac{d_+}{q(m)},\frac{d_-}{q(m)}\right]
\]
are contained in the allowed \(s\)-range.  Their combined length is
\(2(d_+-d_-)/q(m)\), which is at least \(2(d_+-d_-)/q_T^+\).  Therefore
\[
        |\mathcal D_T^\toy(D)|
        \ge \int_{S_T}\frac{2(d_+-d_-)}{q_T^+}\,dm
        =\frac{(d_+-d_-)|I_T|}{q_T^+}.
\]
Dividing by \(|\mathcal D_T|=|I_T|^2\) proves
\eqref{eq:toy-retained-density-exact}.  The estimate
\eqref{eq:toy-retained-density} follows immediately from
\eqref{eq:toy-window-comparability}.  Finally, if
\(|I_T|\ge c_IQ(T)^{-1}\) and \(q_T^-\ge c_q q(T)\), then
\eqref{eq:toy-density-height-hypothesis} follows from
\[
        \frac{d_+}{c_q q(T)}\le \frac{c_I}{4Q(T)},
        \quad\text{i.e.}\quad
        \frac{Q(T)}{q(T)}\le \frac{c_Ic_q}{4d_+},
\]
which holds eventually whenever \(Q(T)=o(q(T))\).
\end{proof}

\begin{lemma}[Signed-separation invariance of the scalar toy channel]
\label{lem:toy-signed-separation-invariance}
For the toy finite blocks of Definition~\ref{def:toy-finite-blocks},
\[
        \widehat\Om_\toy(-s;m,u)=\widehat\Om_\toy(s;m,u)^\top
\]
whenever both sides are defined.  Consequently
\begin{equation}
\label{eq:toy-A-signed-invariance}
        \mathcal A_\toy(\widehat\Om_\toy(-s;m,u))
        =\mathcal A_\toy(\widehat\Om_\toy(s;m,u)).
\end{equation}
\end{lemma}

\begin{proof}
Changing \(s\) to \(-s\) swaps the two centers \(t_-\) and \(t_+\).
The same-point blocks therefore satisfy
\[
        G_{\toy,-}(-s;m,u)=G_{\toy,+}(s;m,u),\qquad
        G_{\toy,+}(-s;m,u)=G_{\toy,-}(s;m,u),
\]
and the mixed block satisfies
\(N_\toy(-s;m,u)=N_\toy(s;m,u)^\top\), by the symmetry of
\(K_\Phi\) and the row/column convention of
Definition~\ref{def:finite-same-point-and-mixed-blocks}.  Positive
definite inverse square roots commute with transposition for symmetric
matrices, so the displayed transpose identity for \(\widehat\Om_\toy\)
follows.  Since
\(\mathcal A_\toy(X)=X_{12}+X_{21}\) is invariant under transpose,
\eqref{eq:toy-A-signed-invariance} follows.
\end{proof}

% =============================================================================
% (2) Insert this corollary at the END of subsec:toy-anomaly-visibility,
%     immediately after the proof of thm:toy-anomaly-visibility.
% =============================================================================

\begin{corollary}[Visibility on the polynomial-weight subdomain]
\label{cor:toy-visibility-polynomial-weight}
Fix \(D\) and the constants \(Q_\toy(D)\), \(u_\toy(D)\), \(c_\toy(D)\)
of Theorem~\ref{thm:toy-anomaly-visibility}.  Assume the height is large
enough that \(q_T^-\ge Q_\toy(D)\), where
\(q_T^-:=\inf_{m\in I_T}q(m)\).  Then for every \(0<|u|\le
u_\toy(D)\) and every \((m,s)\in\mathcal D_T^\toy(D)\),
\begin{equation}
\label{eq:toy-anomaly-visibility-on-subdomain}
        \bigl|\mathcal A_\toy(\widehat\Om_\toy(s;m,u))\bigr|
        \ge \tfrac12 c_\toy(D)u^2.
\end{equation}
If, in addition, the hypotheses of Lemma~\ref{lem:toy-retained-density}
hold, then this lower bound holds on a subset of \(\mathcal D_T\) of
relative density at least the right-hand side of
\eqref{eq:toy-retained-density-exact}, and hence at least
\eqref{eq:toy-retained-density} under
\eqref{eq:toy-window-comparability}.
\end{corollary}

\begin{proof}
Let \((m,s)\in\mathcal D_T^\toy(D)\).  Then
\(d=q(m)|s|\in D\) and \(q_0=q(m)\ge q_T^-\ge Q_\toy(D)\).  If
\(s>0\), Theorem~\ref{thm:toy-anomaly-visibility} gives
\eqref{eq:toy-anomaly-visibility-on-subdomain}.  If \(s<0\), apply the
theorem to \(|s|\) and then use
Lemma~\ref{lem:toy-signed-separation-invariance}.  The density statement
is Lemma~\ref{lem:toy-retained-density}.
\end{proof}

% =============================================================================
% (3) Append this remark at the END of §5, immediately after the
%     existing wip rem:wip-toy-anomaly-uniform-remainder.
% =============================================================================

\begin{remark}[No uniform extension; polynomial-weight scope]
\label{rem:toy-anomaly-no-uniform-extension}
The visibility lower bound \eqref{eq:toy-anomaly-visibility-A} cannot be
extended to all of \(\mathcal D_T\).  At \(s=0\), the same coalescence
argument as Corollary~\ref{cor:omega-coalescence}, applied to the toy
finite blocks, gives \(\widehat\Om_\toy(0;m,u)=I_2\) for every \(u\).
Hence \(\mathcal A_\toy(\widehat\Om_\toy(0;m,u))=0\); no
\(u^2\)-lower bound can hold at the coalescence locus \(d=0\).  The
regimes \(d\to\infty\) and \(d\) approaching an isolated zero of
\(F_\infty\) are excluded by the compactness and avoidance hypotheses on
\(D\) in Definition~\ref{def:toy-retained-subdomain}.  The signed
polynomial-weight subdomain \(\mathcal D_T^\toy(D)\) carries the toy
visibility lower bound by
Corollary~\ref{cor:toy-visibility-polynomial-weight}; its relative
density is polynomial in \(Q\) by Lemma~\ref{lem:toy-retained-density}
and is the only density loss contributed by this replacement of the
impossible all-of-\(\mathcal D_T\) assertion.
\end{remark}

% =============================================================================
% End of referee-revised patch.
% =============================================================================
